\documentclass[10pt,a4paper]{amsart}
\usepackage{amsmath,amssymb,amsthm,graphicx}
\usepackage{verbatim}
\newenvironment{reference}{\comment}{\endcomment}
%\newenvironment{reference}{}{}
\newenvironment{slogan}{\comment}{\endcomment}
\newenvironment{history}{\comment}{\endcomment}

% For commutative diagrams we use Xy-pic
\usepackage[all]{xy}

% We use 2cell for 2-commutative diagrams.
\xyoption{2cell}
% Theorem environments.
%
\theoremstyle{plain}
\newtheorem{theorem}[subsection]{Theorem}
\newtheorem{proposition}[subsection]{Proposition}
\newtheorem{lemma}[subsection]{Lemma}

\theoremstyle{definition}
\newtheorem{definition}[subsection]{Definition}
\newtheorem{example}[subsection]{Example}
\newtheorem{exercise}[subsection]{Exercise}
\newtheorem{situation}[subsection]{Situation}

\theoremstyle{remark}
\newtheorem{remark}[subsection]{Remark}
\newtheorem{remarks}[subsection]{Remarks}

\numberwithin{equation}{subsection}

% Macros
%
\def\lim{\mathop{\mathrm{lim}}\nolimits}
\def\colim{\mathop{\mathrm{colim}}\nolimits}
\def\Spec{\mathop{\mathrm{Spec}}}
\def\Hom{\mathop{\mathrm{Hom}}\nolimits}
\def\Ext{\mathop{\mathrm{Ext}}\nolimits}
\def\SheafHom{\mathop{\mathcal{H}\!\mathit{om}}\nolimits}
\def\SheafExt{\mathop{\mathcal{E}\!\mathit{xt}}\nolimits}
\def\Sch{\mathit{Sch}}
\def\Mor{\mathop{\mathrm{Mor}}\nolimits}
\def\Ob{\mathop{\mathrm{Ob}}\nolimits}
\def\Sh{\mathop{\mathit{Sh}}\nolimits}
\def\NL{\mathop{N\!L}\nolimits}
\def\CH{\mathop{\mathrm{CH}}\nolimits}
\def\proetale{{pro\text{-}\acute{e}tale}}
\def\etale{{\acute{e}tale}}
\def\QCoh{\mathit{QCoh}}
\def\Ker{\mathop{\mathrm{Ker}}}
\def\Im{\mathop{\mathrm{Im}}}
\def\Coker{\mathop{\mathrm{Coker}}}
\def\Coim{\mathop{\mathrm{Coim}}}

% Boxtimes
%
\DeclareMathSymbol{\boxtimes}{\mathbin}{AMSa}{"02}

%
% Macros for moduli stacks/spaces
%
\def\QCohstack{\mathcal{QC}\!\mathit{oh}}
\def\Cohstack{\mathcal{C}\!\mathit{oh}}
\def\Spacesstack{\mathcal{S}\!\mathit{paces}}
\def\Quotfunctor{\mathrm{Quot}}
\def\Hilbfunctor{\mathrm{Hilb}}
\def\Curvesstack{\mathcal{C}\!\mathit{urves}}
\def\Polarizedstack{\mathcal{P}\!\mathit{olarized}}
\def\Complexesstack{\mathcal{C}\!\mathit{omplexes}}
% \Pic is the operator that assigns to X its picard group, usage \Pic(X)
% \Picardstack_{X/B} denotes the Picard stack of X over B
% \Picardfunctor_{X/B} denotes the Picard functor of X over B
\def\Pic{\mathop{\mathrm{Pic}}\nolimits}
\def\Picardstack{\mathcal{P}\!\mathit{ic}}
\def\Picardfunctor{\mathrm{Pic}}
\def\Deformationcategory{\mathcal{D}\!\mathit{ef}}
\usepackage{cleveref}
\crefname{originalproposition}{original proposition}{original propositions}
\Crefname{originalproposition}{Original proposition}{Original propositions}
\makeatletter
\newlabel{derived-proposition-brown}{{0A8G}{source}{}{sourcepointer-1}{}}
\newlabel{derived-proposition-brown@cref}{{[originalproposition][0][]0A8G}{[1][0][]source}{}{}{}}
\makeatother
\crefname{originallemma}{original lemma}{original lemmas}
\Crefname{originallemma}{Original lemma}{Original lemmas}
\makeatletter
\newlabel{perfect-lemma-quasi-coherence-direct-sums}{{08DT}{source}{}{sourcepointer-2}{}}
\newlabel{perfect-lemma-quasi-coherence-direct-sums@cref}{{[originallemma][0][]08DT}{[1][0][]source}{}{}{}}
\makeatother
\crefname{originaltheorem}{original theorem}{original theorems}
\Crefname{originaltheorem}{Original theorem}{Original theorems}
\makeatletter
\newlabel{perfect-theorem-bondal-van-den-Bergh}{{09IS}{source}{}{sourcepointer-3}{}}
\newlabel{perfect-theorem-bondal-van-den-Bergh@cref}{{[originaltheorem][0][]09IS}{[1][0][]source}{}{}{}}
\makeatother
\crefname{originallemma}{original lemma}{original lemmas}
\Crefname{originallemma}{Original lemma}{Original lemmas}
\makeatletter
\newlabel{schemes-lemma-compose-after-separated}{{01KV}{source}{}{sourcepointer-4}{}}
\newlabel{schemes-lemma-compose-after-separated@cref}{{[originallemma][0][]01KV}{[1][0][]source}{}{}{}}
\makeatother
\crefname{originallemma}{original lemma}{original lemmas}
\Crefname{originallemma}{Original lemma}{Original lemmas}
\makeatletter
\newlabel{schemes-lemma-quasi-compact-permanence}{{03GI}{source}{}{sourcepointer-5}{}}
\newlabel{schemes-lemma-quasi-compact-permanence@cref}{{[originallemma][0][]03GI}{[1][0][]source}{}{}{}}
\makeatother
\crefname{originallemma}{original lemma}{original lemmas}
\Crefname{originallemma}{Original lemma}{Original lemmas}
\makeatletter
\newlabel{perfect-lemma-quasi-coherence-pushforward-direct-sums}{{08DZ}{source}{}{sourcepointer-6}{}}
\newlabel{perfect-lemma-quasi-coherence-pushforward-direct-sums@cref}{{[originallemma][0][]08DZ}{[1][0][]source}{}{}{}}
\makeatother
\title{The pushforward along a quasi-compact quasi-separated morphism has a right adjoint}
\author{The Stacks Project Authors}
\date{stacks-project:duality.tex:local-snapshot-20261012; GNU FDL 1.2 or later}
\begin{document}
\maketitle
\noindent\textit{Editorial note.} All delivered mathematical text is the author's own, copied from the pinned Stacks Project chapter file duality.tex; no statement or proof step was rewritten, reordered or removed.\par\medskip
\noindent\textit{Editorial note.} Preamble provenance: the pinned chapter file opens with an input of the shared preamble and carries no preamble text of its own, so the project's single Stacks-wide preamble staged at sources/stacks-duality/preamble.tex is excerpted verbatim under the recipe fixed by LANE\_BRIEF 22.1, as three spans, lines 8-13 (verbatim package and the reference/slogan/history environments), lines 14-18 (Xy-pic setup) and lines 154-217 (the native theorem declarations and mathematical macros). Nothing was added, removed or redefined; the wrapper class and the amsmath/amssymb/amsthm/graphicx base packages are supplied by the assembler.\par\medskip
\noindent\textit{Editorial note.} The counted statement and proof contain no cross reference, no citation and no included file, so the delivered body is self contained.\par\medskip
\noindent\textit{Editorial note.} Canonical location of the counted proposition: Stacks tag 0A9E (sources/stacks-duality/pinned-original-source.tex line 653).\par\medskip
\noindent\textit{Allowed standard prerequisites (editorial).}\begin{itemize}
\item The adjoint functor theorem for compactly generated categories. Derived Categories, Proposition labelled derived-proposition-brown, cited at proof lines 680-682. Supplies the criterion verified in the proof.
\item Direct sums and compact generation of the quasi-coherent derived category. Derived Categories of Schemes, Lemma labelled perfect-lemma-quasi-coherence-direct-sums and Theorem labelled perfect-theorem-bondal-van-den-Bergh, cited at proof lines 684-690. Provide the first two hypotheses.
\item The pushforward commutes with direct sums for a quasi-compact quasi-separated morphism. Derived Categories of Schemes, Lemma labelled perfect-lemma-quasi-coherence-pushforward-direct-sums, cited at proof lines 692-694. Provides the last hypothesis.
\end{itemize}\medskip
\ifdefined\hypertarget\hypertarget{sourcepointer-1}{}\fi
\noindent\textit{Source reference.} Original proposition 0A8G: sources/stacks-chow/tags.txt, source lines 12914--12914. Cross-chapter standard prerequisite located in the project tag list staged in the local snapshot; invoked at proof line 664 of the counted proof. It is retained only as a navigation pointer; every load-bearing step is inside the delivered span.\par\medskip
\ifdefined\hypertarget\hypertarget{sourcepointer-2}{}\fi
\noindent\textit{Source reference.} Original lemma 08DT: sources/stacks-chow/tags.txt, source lines 10648--10648. Cross-chapter standard prerequisite located in the project tag list staged in the local snapshot; invoked at proof line 667 of the counted proof. It is retained only as a navigation pointer; every load-bearing step is inside the delivered span.\par\medskip
\ifdefined\hypertarget\hypertarget{sourcepointer-3}{}\fi
\noindent\textit{Source reference.} Original theorem 09IS: sources/stacks-chow/tags.txt, source lines 12050--12050. Cross-chapter standard prerequisite located in the project tag list staged in the local snapshot; invoked at proof line 670 of the counted proof. It is retained only as a navigation pointer; every load-bearing step is inside the delivered span.\par\medskip
\ifdefined\hypertarget\hypertarget{sourcepointer-4}{}\fi
\noindent\textit{Source reference.} Original lemma 01KV: sources/stacks-chow/tags.txt, source lines 2064--2064. Cross-chapter standard prerequisite located in the project tag list staged in the local snapshot; invoked at proof line 672 of the counted proof. It is retained only as a navigation pointer; every load-bearing step is inside the delivered span.\par\medskip
\ifdefined\hypertarget\hypertarget{sourcepointer-5}{}\fi
\noindent\textit{Source reference.} Original lemma 03GI: sources/stacks-chow/tags.txt, source lines 4529--4529. Cross-chapter standard prerequisite located in the project tag list staged in the local snapshot; invoked at proof line 673 of the counted proof. It is retained only as a navigation pointer; every load-bearing step is inside the delivered span.\par\medskip
\ifdefined\hypertarget\hypertarget{sourcepointer-6}{}\fi
\noindent\textit{Source reference.} Original lemma 08DZ: sources/stacks-chow/tags.txt, source lines 10655--10655. Cross-chapter standard prerequisite located in the project tag list staged in the local snapshot; invoked at proof line 676 of the counted proof. It is retained only as a navigation pointer; every load-bearing step is inside the delivered span.\par\medskip
\begin{lemma}
\label{lemma-twisted-inverse-image}
\begin{reference}
This is almost the same as \cite[Example 4.2]{Neeman-Grothendieck}.
\end{reference}
Let $f : X \to Y$ be a morphism between quasi-separated and quasi-compact
schemes. The functor $Rf_* : D_\QCoh(\mathcal{O}_X) \to D_\QCoh(\mathcal{O}_Y)$
has a right adjoint.
\end{lemma}

\begin{proof}
We will prove a right adjoint exists by verifying the hypotheses of
Derived Categories, Proposition \ref{derived-proposition-brown}.
First off, the category $D_\QCoh(\mathcal{O}_X)$ has direct sums, see
Derived Categories of Schemes, Lemma
\ref{perfect-lemma-quasi-coherence-direct-sums}.
The category $D_\QCoh(\mathcal{O}_X)$ is compactly generated by
Derived Categories of Schemes, Theorem
\ref{perfect-theorem-bondal-van-den-Bergh}.
Since $X$ and $Y$ are quasi-compact and quasi-separated, so is $f$, see
Schemes, Lemmas \ref{schemes-lemma-compose-after-separated} and
\ref{schemes-lemma-quasi-compact-permanence}.
Hence the functor $Rf_*$ commutes with direct sums, see
Derived Categories of Schemes, Lemma
\ref{perfect-lemma-quasi-coherence-pushforward-direct-sums}.
This finishes the proof.
\end{proof}

\end{document}
